
\chapter{Introduction}
\begin{refsection}


\Acp{BH} are among the most extreme and enigmatic objects in nature. Yet despite over a century of theoretical work and multiple decades of increasingly detailed and sophisticated observation, the astrophysical formation processes remain elusive. 


There is a long tradition of theoretical work on \acp{BH}. The first recognizable ideas can be traced back to \citet{Michell:1784xqa} who speculated about ``dark stars,'' stars so compact that the escape velocity from the surface exceed the speed of light. Newtonian corpuscles of light could therefore not escape the gravitational potential of the compact object: Michell's dark star was impossible to observe with traditional telescopes. 

Several scientific revolutions later, Michell's speculation was modernized by \citet{Schwarzschild:1916sss} as a spherically symmetric solution to Einstein's new gravitational general theory of relativity (GR) \citep{Einstein:1916vd}. Schwarzschild's solution included a causal horizon beyond which signals---light, in particular---cannot escape. Schwarzschild \acp{BH} cannot be directly observed with light. 

Early theoretical work on \acp{BH} rarely entertained that \ac{BH} might be physical objects that exist in nature. Over the next sixty years, however, the theoretical case for astrophysical \acp{BH} strengthened with the emergence of singularity theorems, most notably due to Penrose \citep{}, showing that collapse is not a byproduct of idealized symmetry, but a robust prediction of general relativity. When coupled with the existence of a maximum mass for degenerate matter---the Tolman-Oppenheimer-Volkoff limit \citep{}---a \ac{BH} became the only theoretical refuge for sufficiently massive stellar remnants.



\printbibliography[heading=subbibliography]
\end{refsection}